\chapter{Conclusão}\label{sec:conclusao}

\begin{dica}[Dica de metodologia científica: o que escrever na Conclusão]
A conclusão \emph{responde} aos objetivos: retome o objetivo geral e, um a um, os
específicos, dizendo o que foi alcançado e apontando a evidência (seção, resultado,
artefato). Em seguida: (i) as \textbf{contribuições} do trabalho; (ii) as
\textbf{limitações} (o que não foi feito ou não pode ser generalizado); e (iii) os
\textbf{trabalhos futuros}, de forma concreta. Não introduza dados, figuras ou citações
novas, não repita a introdução nem termine com frases genéricas. O texto deve ser
compreensível para quem leu apenas a introdução e a conclusão.
\end{dica}

Apresente as conclusões correspondentes aos objetivos ou hipóteses, as limitações e
os trabalhos futuros.
